\subsection{События автопилота}

Таймер --- это событие «прошло время». Автопилот умеет сообщать и о других событиях: моторы запущены, взлёт завершён, точка достигнута, дрон сел, батарея разряжена, удар. Для них есть одна функция:

\begin{codebox}[Обработчик событий]
\begin{verbatim}
function callback(event)
  if event == Ev.COPTER_LANDED then
    -- что сделать после посадки
  end
end
\end{verbatim}
\end{codebox}

Мы её только определяем. Вызывает её автопилот --- каждый раз, когда что-то произошло, --- и передаёт в \verb|event| код события. Коды хранятся в таблице \verb|Ev|; сравнивать нужно именно с \verb|Ev.ИМЯ|, а не с числом.

\begin{center}
\small
\begin{tabularx}{\linewidth}{@{}>{\ttfamily}l X@{}}
\toprule
\normalfont\textbf{Событие} & \textbf{Когда приходит} \\
\midrule
Ev.ENGINES\_STARTED & моторы запущены после предстартовой подготовки \\
Ev.TAKEOFF\_COMPLETE & взлёт завершён, дрон висит на высоте взлёта \\
Ev.POINT\_REACHED & достигнута точка, заданная \texttt{ap.goToLocalPoint} \\
Ev.COPTER\_LANDED & дрон коснулся земли после посадки \\
Ev.LOW\_VOLTAGE1, Ev.LOW\_VOLTAGE2 & батарея разряжена: первый порог --- предупреждение, второй --- пора садиться \\
Ev.SHOCK & удар или столкновение \\
\bottomrule
\end{tabularx}
\end{center}

\subsubsection*{Реакция на посадку}

\begin{codefile}{lua/05_leds/leds_events_landed.lua}
\end{codefile}

Линейка горит белым всё время работы и гаснет, когда автопилот сообщит о посадке. В симуляторе проверить это можно так: добавьте в начало программы взлёт и посадку из главы~\ref{sec:missions} --- или запустите готовый пример \texttt{leds\_mission\_stages.lua} из следующего раздела.

\subsubsection*{Реакция на разряд батареи}

\begin{codefile}{lua/05_leds/leds_events_low_voltage.lua}
\end{codefile}

\paragraph{Что происходит.} В штатном режиме таймер мигает зелёным. Событие \verb|LOW_VOLTAGE2| останавливает таймер и через секунду заливает линейку красным. Пауза нужна, потому что уже запланированный тик таймера может сработать ещё раз сразу после \verb|stop()| --- так устроен автопилот Пионера.

Короткий вариант того же приёма, без таймера мигания. Останавливать здесь нечего, поэтому секундная пауза не обязательна --- это просто отложенная реакция на событие:

\begin{codefile}{lua/05_leds/leds_timers_callback_delayed.lua}
\end{codefile}

\begin{rulebox}[callback короткий и быстрый]
Внутри \verb|callback| --- только смена состояния, цвет и команда автопилоту. Никаких \verb|sleep| и длинных циклов: пока обработчик работает, следующие события ждут.
\end{rulebox}
